\section*{Erd\H{o}s Problem 421}

\paragraph{Statement and scope.}
Does there exist \(A=\{a_1<a_2<\cdots\}\subset\mathbb N\) of natural
density one such that all products of nonempty consecutive blocks are
distinct? The current page records an affirmative solution.
This attempt reconstructs the prime-gap construction in Kyle Pratt's
exposition and proves its product-distinctness. The density-one estimate
is unresolved here and is isolated precisely below.

\paragraph{A construction that retains every prime.}
Call a finite increasing set product-distinct if two nonempty consecutive
blocks with equal products have the same endpoints. Set \(A_2=\{2\}\).
For consecutive primes \(p<q\), form
\[
 B=A_p\cup\{p+1,\ldots,q\}.
\]
If \(B\) is product-distinct, set \(A_q=B\).
Otherwise set \(A_q=A_p\cup\{q\}\), and call \((p,q)\) rejected.
The test is finite: list all consecutive products and check for repeats.
Define \(A=\bigcup_p A_p\).

\paragraph{Proof that no equal consecutive products occur.}
Suppose \(A_p\) is product-distinct. Acceptance preserves this by
definition. For rejection, every old element is less than the prime
\(q\), so no product of old elements is divisible by \(q\).
Consequently a collision in \(A_p\cup\{q\}\) would have to involve
two blocks both containing \(q\), or two blocks both avoiding it.
The second case is excluded by induction. In the first case both blocks
end at \(q\); after canceling their common terminal factors, different
starting points would give a nonempty product of integers at least
two equal to \(1\), which is impossible. Thus induction applies.

The sets \(A_p\) are nested and future stages only append larger
elements. Every pair of finite blocks of \(A\) is therefore already
present as the same pair of blocks in some \(A_p\). They cannot collide.
Also every prime is retained, so \(A\) is infinite.
Initial stages are
\[
\begin{split}
 A_5&=\{2,3,4,5\},\\
 A_7&=\{2,3,4,5,7\},\\
 A_{13}&=\{2,3,4,5,7,8,9,10,11,13\},\\
 A_{17}&=A_{13}\cup\{17\}.
\end{split}
\]
Here \(6\) is rejected because \(2\cdot3=6\), \(12\) because
\(3\cdot4=12\), and the interval \((13,17)\) because
\(3\cdot4\cdot5\cdot7\cdot8=14\cdot15\cdot16\).
These stages and the absence of other collisions were checked exactly.

\paragraph{The exact density estimate still needed.}
The omitted integers consist of \(1\) and the interiors of rejected
prime gaps. For every real \(X\geq1\),
\[
 X-\lfloor X\rfloor+\bigl|[1,X]\cap(\mathbb N\setminus A)\bigr|
 =X-|A\cap[1,X]|.
\]
In particular, density one is equivalent to
\[
 \sum_{\substack{p<q\ {\rm consecutive\ primes}\\(p,q)\ {\rm rejected}}}
 \bigl|(p,q)\cap[1,X]\cap\mathbb Z\bigr|=o(X).
 \tag{1}
\]
The partial final gap is included, so this criterion applies to arbitrary
\(X\), not merely prime endpoints.

One can separate (1) into gaps of length at most \(p^\theta\) and
longer gaps, with \(\theta=5/72\). It suffices to prove that each of the
two corresponding weighted sums is \(o(X)\). Counting rejected gaps
without their lengths would not suffice.
Pratt's proof supplies the necessary estimates through results on
primes in almost all short intervals and bounds for integer points
on algebraic curves, together with a recursive analysis of rejection.
Those estimates are the missing work here; the elementary construction
alone does not establish density one.

\paragraph{Sources and attribution.}
Kyle Pratt, \emph{A digested proof of Erd\H{o}s Problem 421},
\url{https://www.erdosproblems.com/static/421-Pratt.pdf},
especially Section 2. Pratt credits earlier AI-assisted work by
Chojecki and Sneiderman and discusses a different proof by Sharma.
Current statement: \url{https://www.erdosproblems.com/421},
checked 24 September 2026. This is a partial reconstruction with
explicit attribution, not a claim of a new solution.

\paragraph{Completion estimate.}
The infinite product-distinct construction is proved; the decisive density estimates are missing.
\textbf{COMPLETION: 40\%}
